% !TeX root = ../../Book4.tex
% 纲要标题：### 11.3 零调对象
% 纲要标签：b4:sec:1003
% 纲要位置：计划纲要.md，第 3478 行。

\section{零调对象}
\label{b4:sec:1003}

计算一个指定函子时，消解项不必全部内射，但必须对该函子没有正次导出贡献。仅有原列正合并不足够，逐项的零调条件还需证明。

% 纲要内容：
%
% * 相对于函子的零调性
% * 使用零调消解计算导出函子
% * 不能随意用任意正合复形替代消解
%

\subsection{相对于函子的零调性}
\label{b4:subsec:1003:01}

内射消解对每个左正合函子都可用，但有时它的各项过大，实际计算并不方便。若只计算一个指定函子 \(F\)，可以使用对这个函子没有高次导出贡献的对象。

\begin{definition}\label{b4:def:functor-acyclic}
设 \(\mathcal A,\mathcal B\) 为交换范畴，\(\mathcal A\) 有足够多内射对象，\(F:\mathcal A\to\mathcal B\) 为左正合加性函子。若对象 \(Q\in\mathcal A\) 满足 \(R^iF(Q)=0\) 对所有整数 \(i>0\) 成立，则称 \(Q\) 为 \(F\)-\term[lingtiao-duixiang]{零调对象}[acyclic object]，也称 \(F\)-无上同调对象\termidx[wushangtongdiaoduixiang]{无上同调对象}。给定 \(\mathcal A\) 中的对象 \(M\) 及非负增广上链复形
\[
 0\longrightarrow M\longrightarrow A^0\longrightarrow A^1\longrightarrow\cdots
\]
若此序列正合，且每个 \(A^j\)（\(j\ge0\)）都为 \(F\)-零调对象，则称为 \(M\) 的 \(F\)-\term[lingtiao-xiaojie]{零调消解}[acyclic resolution]。
\end{definition}

这里的零调性描述一个对象相对于指定函子 \(F\) 的高次导出行为，只要求正次数的导出函子消失，不要求 \(F(Q)=0\)。\cref{b4:def:acyclic-complex}中的零调复形则描述整条复形内部的同调，要求所有次数的同调对象为零。零调消解要求各项对 \(F\) 零调；去掉增广项后的复形有 \(H^0(A^\bullet)\cong M\)，而施加 \(F\) 后的正次数上同调也可能非零。

\ProseParagraphBreak
内射对象对每个左正合加性函子都是零调对象；对于固定的函子，零调对象却未必内射。例如取 \(F=\operatorname{Hom}_{\mathbb Z}(\mathbb Z/2\mathbb Z,-)\)。对任意交换群 \(Q\)，长度一自由消解 \(0\to\mathbb Z\xrightarrow{\times2}\mathbb Z\to\mathbb Z/2\mathbb Z\to0\) 给出
\[
 R^1F(Q)\cong Q/2Q,\qquad R^iF(Q)=0\quad(i>1).
\]
因此 \(\mathbb Z/3\mathbb Z\) 对 \(F\) 零调，却因不可除而不是内射交换群；不过此时 \(F(\mathbb Z/3\mathbb Z)=0\)。若改取
\[
 Q=(\mathbb Q/\mathbb Z)\oplus\mathbb Z/3\mathbb Z,
\]
则 \(2Q=Q\)，但 \(3Q\ne Q\)，所以它仍对 \(F\) 零调而非内射。同时 \(F(Q)\cong Q[2]\cong\mathbb Z/2\mathbb Z\ne0\)，明确说明原函子的取值不必消失。

\begin{proposition}\label{b4:prop:acyclic-closure}
设 \(\mathcal A,\mathcal B\) 为交换范畴，\(\mathcal A\) 有足够内射对象，\(F:\mathcal A\to\mathcal B\) 为左正合加性函子，且 \(0\to A\to B\to C\to0\) 为 \(\mathcal A\) 中的短正合列。若 \(A,C\) 都 \(F\)-零调，则 \(B\) 也如此；若 \(A,B\) 都 \(F\)-零调，则 \(C\) 也如此。若 \(B,C\) 都 \(F\)-零调，则 \(A\) 零调当且仅当 \(F(B)\to F(C)\) 为满态射。
\end{proposition}
\begin{proof}
前两项分别由长正合列中的
\(R^iF(A)\to R^iF(B)\to R^iF(C)\) 及
\(R^iF(B)\to R^iF(C)\to R^{i+1}F(A)\) 得到，其中 \(i\ge1\)。当 \(B,C\) 都零调时，片段 \(R^{i-1}F(C)\to R^iF(A)\to R^iF(B)\) 已使 \(R^iF(A)=0\) 对 \(i\ge2\) 成立，但低次部分仍为
\[
 0\longrightarrow F(A)\longrightarrow F(B)\longrightarrow F(C)
 \longrightarrow R^1F(A)\longrightarrow0.
\]
所以唯一可能剩下的正次导出是 \(R^1F(A)\cong\operatorname{coker}(F(B)\to F(C))\)。它消失恰好要求 \(F(B)\to F(C)\) 满，这就是第三种情形多出的条件。
\end{proof}

\subsection{零调消解与导出函子的计算}
\label{b4:subsec:1003:02}

\begin{theorem}\label{b4:thm:acyclic-resolution}
设 \(\mathcal A,\mathcal B\) 为交换范畴，\(\mathcal A\) 有足够内射对象，\(F:\mathcal A\to\mathcal B\) 为左正合加性函子。若 \(0\to M\to A^0\to A^1\to\cdots\) 是 \(F\)-零调消解，则对每个 \(n\ge0\)，有典范同构
\[
 R^nF(M)\simeq H^n(F(A^\bullet)).
\]
此同构对增广消解之间的态射自然。
\end{theorem}
\begin{proof}
令 \(K^0=M\)，从增广嵌入 \(\iota^0:M\hookrightarrow A^0\) 起，逐次取 \(K^{j+1}=\operatorname{coker}\iota^j\)。消解正合性把这个商对象识别为 \(\operatorname{im}(\mathrm{d}_A^j)\)，并给出下一项中的嵌入 \(\iota^{j+1}:K^{j+1}\hookrightarrow A^{j+1}\)。这样前两步是
\[
\begin{gathered}
 0\longrightarrow M\longrightarrow A^0\longrightarrow K^1\longrightarrow0,\\
 0\longrightarrow K^1\longrightarrow A^1\longrightarrow K^2\longrightarrow0,
\end{gathered}
\]
一般的一步写成
\[
 0\longrightarrow K^j\xrightarrow{\iota^j} A^j\xrightarrow{p^j} K^{j+1}\longrightarrow0
 \qquad(j\ge0).
\]
先看一次上同调。微分 \(\mathrm{d}_A^1=\iota^2p^1\)，左正合性使 \(F(\iota^2)\) 单，并把 \(\ker F(p^1)\) 识别为 \(F(K^1)\)。因而施加 \(F\) 后的一次余圈对象是 \(F(K^1)\)，一次余边来自 \(F(A^0)\to F(K^1)\)。第一条短正合列的长正合列又因 \(A^0\) 零调而给出 \(F(A^0)\to F(K^1)\to R^1F(M)\to0\)，所以
\[
 H^1(F(A^\bullet))
 \cong\operatorname{coker}\bigl(F(A^0)\to F(K^1)\bigr)
 \cong R^1F(M).
\]
一般次数的计算把同一机制用于后面的短正合列。长正合列和 \(A^j\) 的零调性给出
\[
\begin{gathered}
 0\to F(K^j)\to F(A^j)\to F(K^{j+1})\to R^1F(K^j)\to0,\\
 R^iF(K^{j+1})\simeq R^{i+1}F(K^j)\qquad(i\ge1).
\end{gathered}
\]
消解的微分分解为
\[
 \mathrm{d}_A^n:\quad A^n\xrightarrow{p^n}K^{n+1}
 \xrightarrow{\iota^{n+1}}A^{n+1}.
\]
对第 \(n+1\) 步的短正合列用左正合性，得到 \(F(\iota^{n+1})\) 仍为单态射，所以复合它不改变 \(F(p^n)\) 的核。再对第 \(n\) 步用左正合性，\(F(\iota^n)\) 把 \(F(K^n)\) 识别为这个核。因此
\[
\begin{aligned}
 \ker\bigl(F(\mathrm{d}_A^n)\bigr)
 &=\ker\bigl(F(p^n)\bigr)\\
 &=F(K^n),
\end{aligned}
\]
即 \(F(A^\bullet)\) 在次数 \(n\) 的余圈对象为 \(F(K^n)\)。当 \(n\ge1\) 时，同调就是
\[
 \operatorname{coker}\bigl(F(A^{n-1})\to F(K^n)\bigr)
 \simeq R^1F(K^{n-1})\simeq R^nF(M).
\]
次数零则由左正合性直接等于 \(F(M)\)。这些同构只使用给定的正合列和连接同态，未选分裂，因而对消解态射自然。

为说明这里的“典范”也与最初采用内射消解的定义一致，取 \(M\to I^\bullet\)。下面只使用目标 \(I^j\) 的内射性，源 \(A^j\) 不必内射：先把 \(M\to I^0\) 沿 \(M\hookrightarrow A^0\) 延拓为 \(f^0:A^0\to I^0\)。若已构造到 \(f^j\)，则 \(\mathrm{d}_I^jf^j\) 在 \(K^j\) 上为零；在零次这是与增广相容，在正次这是已有链映射等式及 \(\mathrm{d}_I^j\mathrm{d}_I^{j-1}=0\) 的结果。于是该复合经余核 \(p^j\) 分解为 \(K^{j+1}\to I^{j+1}\)，再由 \(I^{j+1}\) 内射沿 \(\iota^{j+1}\) 延拓为 \(f^{j+1}\)。逐次得到增广上链映射 \(f^\bullet:A^\bullet\to I^\bullet\)：
\[
\begin{tikzcd}[column sep=2.2em,row sep=large]
0\arrow[r]&M\arrow[r]\arrow[d,equal]&A^0\arrow[r]\arrow[d,"f^0"]&A^1\arrow[r]\arrow[d,"f^1"]&\cdots\\
0\arrow[r]&M\arrow[r]&I^0\arrow[r]&I^1\arrow[r]&\cdots
\end{tikzcd}
\]
这里的比较映射通过延拓选取，其同伦唯一性由\cref{b4:prop:comparison-homotopy-unique}的内射版本保证。具体地，设 \(f^\bullet,g^\bullet:A^\bullet\to I^\bullet\) 都延拓 \(1_M\)，约定 \(I^{-1}=0\)、\(h^0=0\)。由于 \(f^0-g^0\) 在 \(M\) 上为零，它经 \(p^0:A^0\twoheadrightarrow K^1\) 分解为 \(K^1\to I^0\)。利用 \(I^0\) 的内射性，把后者沿 \(\iota^1:K^1\hookrightarrow A^1\) 延拓为 \(h^1:A^1\to I^0\)，于是 \(f^0-g^0=h^1\mathrm{d}_A^0\)。

假定已构造到 \(h^j\)（\(j\ge1\)），且低于 \(j\) 次的同伦等式成立。已选的 \(h^j\) 解释了由前一次数传来的差别，在次数 \(j\) 尚未消去的误差是
\[
 e^j=f^j-g^j-\mathrm{d}_I^{j-1}h^j.
\]
链映射等式及前一次的同伦等式给出
\[
\begin{aligned}
 e^j\mathrm{d}_A^{j-1}
 &=\mathrm{d}_I^{j-1}
   \bigl(f^{j-1}-g^{j-1}-h^j\mathrm{d}_A^{j-1}\bigr)\\
 &=\mathrm{d}_I^{j-1}\mathrm{d}_I^{j-2}h^{j-1}=0.
\end{aligned}
\]
这个等式说明误差在前一微分的像上为零。由于 \(\mathrm{d}_A^{j-1}=\iota^jp^{j-1}\) 且 \(p^{j-1}\) 为满态射，\(e^j\iota^j=0\)，故 \(e^j\) 经余核 \(p^j\) 下降为 \(\bar e^j:K^{j+1}\to I^j\)。再由 \(I^j\) 内射，把 \(\bar e^j\) 沿 \(\iota^{j+1}\) 延拓为 \(h^{j+1}:A^{j+1}\to I^j\)。这样逐次得到
\[
 f^j-g^j=\mathrm{d}_I^{j-1}h^j+h^{j+1}\mathrm{d}_A^j
 \qquad(j\ge0),
\]
即 \(f-g=\mathrm{d}_Ih+h\mathrm{d}_A\)。施加加性函子后，这个等式仍成立，所以两种比较映射诱导同一个同调映射。上面的短正合列在比较映射下构成交换图，连接同态自然性说明其诱导映射正是已经构造的同构。因此不同零调消解都通过给定的 \(R^nF(M)\) 典范比较，而不要求它们之间任意方向都有链映射。
\end{proof}

逐项零调并不使 \(F(A^\bullet)\) 成为零调复形。它只使各 \(A^j\) 的正次导出消失，却未要求 \(F(A^{j-1})\to F(K^j)\) 满。上面计算的余核正是这份满射性失败；沿短正合列平移后，它记录的是 \(R^jF(M)\)。

\ProseParagraphBreak
对有足够投射对象的交换范畴及右正合加性函子，对定义域和值域同时取相反范畴，就得到左导出版本：若 \(Q_\bullet\to M\) 正合，且 \(L_iF(Q_j)=0\) 对所有 \(i>0,j\ge0\) 成立，则 \(L_nF(M)\simeq H_n(F(Q_\bullet))\)。在相反范畴中，核与余核互换，内射延拓成为投射提升；再把上链次数改记为链次数，便得到这些正合列与同构，而不需额外假设。具体地，设 \(R\) 为含幺环，\(N\) 为幺性左 \(R\) 模。在幺性右 \(R\) 模范畴中取 \(F=-\otimes_RN\)，平坦右模 \(Q\) 满足 \(L_iF(Q)=\operatorname{Tor}_i^R(Q,N)=0\)（\(i>0\)）。因此任一平坦消解 \(\cdots\to Q_1\to Q_0\to M\to0\) 都给出
\[
 \operatorname{Tor}_n^R(M,N)\cong H_n(Q_\bullet\otimes_RN)
 \qquad(n\ge0),
\]
而不必要求各 \(Q_j\) 投射。

\subsection{正合复形作为消解的条件}
\label{b4:subsec:1003:03}

应用定理时，要检查增广列正合、各项对指定函子零调，并且消解从次数零沿一个方向展开。前两项给出连接同态的平移，最后一项保证每个固定的 \(n\ge2\) 经过有限次就回到 \(M\)：
\[
 R^1F(K^{n-1})\cong R^2F(K^{n-2})\cong\cdots\cong R^nF(K^0)
 =R^nF(M).
\]
一条向两端无限延伸的正合复形若没有这样的增广起点，即使局部短正合列允许平移，也不能据此完成与指定对象导出的比较。

\begin{example}\label{b4:exm:arbitrary-resolution-fails}
取 \(F=\operatorname{Hom}_{\mathbb Z}(\mathbb Z/2\mathbb Z,-)\) 和 \(M=\mathbb Z/2\mathbb Z\)。正合列
\[
 0\longrightarrow M\longrightarrow M\oplus\mathbb Z
 \longrightarrow\mathbb Z\longrightarrow0
\]
用自然包含和投影，确实是一个长度一的右消解。但施加 \(F\) 后，非增广上链复形为 \(\mathbb Z/2\mathbb Z\to0\)，一次上同调为零，而 \(R^1F(M)=\operatorname{Ext}_{\mathbb Z}^1(\mathbb Z/2\mathbb Z,\mathbb Z/2\mathbb Z)=\mathbb Z/2\mathbb Z\)。第零项的零调条件已经失败，因为
\[
 R^1F(M\oplus\mathbb Z)
 \cong (M\oplus\mathbb Z)/2(M\oplus\mathbb Z)
 \cong (\mathbb Z/2\mathbb Z)^{\oplus2}\ne0.
\]
末项 \(\mathbb Z\) 也有 \(R^1F(\mathbb Z)\cong\mathbb Z/2\mathbb Z\ne0\)。原列虽正合，却不是 \(F\)-零调消解，不能用它计算导出函子。
\end{example}

对同一个 \(M=\mathbb Z/2\mathbb Z\)，取
\[
 0\longrightarrow\mathbb Z/2\mathbb Z\xrightarrow{\eta}\mathbb Q/\mathbb Z
 \xrightarrow{\times2}\mathbb Q/\mathbb Z\longrightarrow0,
 \qquad \eta(\bar1)=\dfrac12+\mathbb Z.
\]
乘二映射为满态射，其核恰为 \(\eta(\mathbb Z/2\mathbb Z)\)，而 \(\mathbb Q/\mathbb Z\) 是可除群，因而此列为内射消解。\(F(\mathbb Q/\mathbb Z)\) 由 \(\mathbb Q/\mathbb Z\) 中被 \(2\) 零化的元素确定，同构于 \(\mathbb Z/2\mathbb Z\)；乘二在这些元素上为零，所以施加同一 \(F\) 后，非增广上链复形为
\[
 \mathbb Z/2\mathbb Z\xrightarrow{0}\mathbb Z/2\mathbb Z
 \qquad\text{（次数 }0,1\text{）}.
\]
于是
\[
 R^0F(\mathbb Z/2\mathbb Z)\simeq\mathbb Z/2\mathbb Z,
 \qquad R^1F(\mathbb Z/2\mathbb Z)\simeq\mathbb Z/2\mathbb Z.
\]
这两个正合消解针对同一个对象；前一个含有非零调项，后一个逐项内射，因此后者能够计算该对象的导出函子。

\ProseParagraphBreak
零调对象及消解计算可参阅 Weibel\cite[F-Acyclic Objects 2.4.3, Exercise 2.4.3, Section 2.5]{Weibel1994HomologicalAlgebra}。该书在左导出部分提出此类对象，并以维数平移给出练习；本节把当前所需的右导出版本作为正文定理证明。
